\section*{Methods}
\label{sec:method}

%%% NOTES. (journal instructions)
%%%
%%% The Methods should include detailed text describing any steps or procedures used in producing the data, including full descriptions of the experimental design, data acquisition assays, and any computational processing (e.g. normalization, image feature extraction). See the detailed section in our submission guidelines for advice on writing a transparent and reproducible methods section. Related methods should be grouped under corresponding subheadings where possible, and methods should be described in enough detail to allow other researchers to interpret and repeat, if required, the full study. Specific data outputs should be explicitly referenced via data citation (see Data Records and Citing Data, below). Authors should cite previous descriptions of the methods under use, but ideally the method descriptions should be complete enough for others to understand and reproduce the methods and processing steps without referring to associated publications. There is no limit to the length of the Methods section. Subheadings should not be numbered.


To compile a dataset on text reuse in scientific writing, three things are required: a large collection of scientific publications; detailed metadata for each of these publications; and a scalable, yet accurate method for detecting reuse between the compiled texts. This section formally introduces all three of these components. Our pipeline and its individual steps are depicted in Figure~\ref{fig:pipeline} for reference. Note that text reuse detection is typically implemented as a bipartite process of source retrieval to identify candidate pairs of documents, followed by text alignment to identify reused portions of text between both candidate documents \cite{stein:2007d}.


\subsection*{Publication Acquisition, Preprocessing and Metadata}\label{sub:data}

Before the detection of scientific text reuse at scale can commence, a large amount of scientific publications and detailed metadata has to be collected. We assemble such data in the five steps document selection, plain text extraction, text preprocessing, meta data acquisition, and meta data standardization.

For raw document selection, we make use of the CORE dataset \cite{knoth:2012}, one of the largest collections of open access academic publications, obtained from more than 12,000 data providers. First, we identify the 6,015,512 unique open-access DOIs in the CORE dataset dump from March~1st,~2018. Since the plain texts extracted from the publications' PDF files as supplied by the CORE data are of highly varying quality, and since no structural annotations (such as markup for citations, in-text references, section annotations) are available, we opted to obtain the original PDF files of the identified open access DOIs from various openly accessible repositories.

Plain text extraction was then repeated on the acquired PDF files using the standardized state-of-the-art toolchain GROBID \cite{lopez:2015}. The extracted texts are preprocessed as follows: For each document, abstract and main text paragraphs are extracted as indicated by the content tags provided by GROBID; in-text references, tables, and figures are omitted. We choose to further drop the bibliographical data, special characters, and numerical characters to reduce the amount of false positives for the subsequent text alignment step. Since the retrieval models underlying the text alignment detection compute the word overlap between document passages, standardized or meta text fragments such as citations, numbers, author names, affiliations, and references exponentially increase the set of passages matched between documents that otherwise may not share text. A minimum of 1,000 and a maximum of 60,000 whitespace-separated words is imposed as an effective heuristic to filter common plain text extraction errors. Overall, we obtained and extracted clean plain text for 4,267,166 documents~(70\% of the original CORE dataset).

Given this set, we cross-check and supplement the metadata provided by CORE with additional data from the Microsoft Open Academic Graph \cite{tang:2008,sinha:2015}, which features fields-of-study annotations for a wide set of publications. Since the annotated disciplines do not follow a hierarchical scheme, and since they are of differing granularity per publication (e.g., ``humanities'' as a whole vs.\ ``chemical solid-state research'' as a subfield of chemistry), we manually map the classification found in the Microsoft Open Academic Graph to the standardized, hierarchical \textit{DFG Classification of Scientific Disciplines, Research Areas, Review Boards and Subject Areas} \cite{dfg:2016}.\footnote{We opt to replace the term \textit{review board} as used in the DFG classification with the more conventional denomination \textit{field of study}.}



\subsection*{Text Reuse Detection in Large Document Collections}

Given the large collection of plain texts, a text reuse detection approach is applied  to identify reused passages between pairs of texts. At heart, the detection of text reuse boils down to the task of {\em text alignment} \cite{potthast:2013e}: Given two documents~$d$ and~$d'$, determine the set $R_{d,d'}$ of all pairs of reused text passages $(t, t')$, where $t\subseteq d$ and $t' \subseteq d'$. 

When given a set of documents~$D$, the computational complexity of detecting the set $R_D$ of all pairs of reused text passages among all pairs of documents, $\{(d, d')\mid d, d'\in D\}$, grows quadratically with the document count, {$|D|$}. If an exhaustive comparison exceeds the available computing capacity, a source retrieval step has to be carried out first \cite{hagen:2015e,hagen:2017e}: Given a document~$d$ and a document collection~$D$, a subset $D_d\subseteq D$ of candidate documents is retrieved containing those documents $d'\in D_d$ whose likelihood of sharing a reused text passage with~$d$ is sufficiently high. The size $|D_d|$ for a given document~$d$ varies dependent on how many documents in~$D$ have a sufficiently high reuse likelihood, but is expected to be $|D_d| \ll |D|$. Note that the estimation of a reuse likelihood for all pairs $\{(d, d')\mid d, d'\in D\}$ can be operationalized with linear time complexity in~$|D|$, e.g., by using locality-sensitive hashing \cite{stein:2019c}.

Source retrieval introduces an error in terms of lower recall, i.e., missed pairs of reused passages that would have been detected by an exhaustive text alignment of every document pair. The goal of source retrieval therefore is to closely approximate the recall of an exhaustive comparison while minimizing the sum of the sizes of retrieved candidate documents $\sum_{d\in D} |D_d|$, i.e., the computing capacity actually expended. Altogether, to detect reused text passages in~$D$, we retrieve for each $d \in D$ the set $D_d$ of candidate documents and compute for each $d'\in D_d$ the set of reused text passages $R_{d,d'}$. The union of all detected passages is used as an approximation of~$R_D$.



\subsection*{Source Retrieval}

The source retrieval component of our pipeline, i.e., the computation of $D_d$ for a given $d$, is operationalized by treating text reuse as a locally bounded phenomenon. Candidate pairs are presumed to be identifiable by comparing text passages between two documents. If the similarity is sufficiently high for at least one combination of passages between two documents $d, d'$, then $d'$ is considered a candidate for $d$. To achieve high scalability, we apply a locality-sensitive hash function~$h$ to each passage~$t$ in a document~$d$, thus representing each document as a set of hash values~$h(t)$. The similarity between two passages $t,t'$ is then approximated by the extent of overlap between their hash sets, $|h(t)\cap h(t')|$. Thus, a document $d'$ is considered a candidate source of text reuse for a document $d$, if at least one of their passage-level hash sets intersects \cite{stein:2019c}: $\exists t \subseteq d \ \exists t' \subseteq d' : h(t) \cap h(t')\neq\emptyset$. 

In our pipeline, the documents are first divided into consecutive passages of $n$~terms, which in turn are each embedded using a bag-of-words representation, i.e., as their term occurrence vector. To find matching passages between two documents, MinHash \cite{broder:97} is chosen as hashing scheme for~$h$. For each passage vector, multiple individual hash functions are applied to each element, and the minimum hash value of each pass is saved, yielding the set minimum hashes for a passage. It can be shown that, for~$m$ individual hashes per passage, two texts with a Jaccard similarity of at least~$m^{-1}$ are guaranteed to produce a hash collision between their hash sets, rendering them a suitable approximation for the lower bound of document similarity \cite{broder:97}.

Hash-based source retrieval allows for a significant reduction of the required computation time, since the hash-based approximation has a linear time complexity with respect to~$|D|$, as opposed to the quadratic complexity of vector comparisons \cite{stein:2019c}. As a result, our source retrieval computation time could be fitted into the allotted budget of two months of computing time on a 130-node Apache Spark cluster, with 12~CPUs and 196~GB RAM per node. 

Besides its efficiency, the outlined approach provides us with a second highly beneficial characteristic---the source retrieval step depends on only two parameters: the passage length~$n$, which can be chosen according to computational constraints, and the number of hashes~$m$, which is chosen according to the required minimum similarity separating two passages. Applying this scheme to all passages in all documents, we detect all document pairs that are locally similar with at least~$m^{-1}$ Jaccard similarity, i.e., which have a word overlap in their bag-of-words passage representations. By choosing this bound lower than the minimum detection threshold of the subsequent text alignment step, the search space is pruned without a loss in accuracy. Thus, our hash-based similarity search primarily impacts the precision, but not the recall of the source retrieval step with respect to the subsequent text alignment step.



\subsection*{Text Alignment}\label{sub:alignment}

Given a set of candidates $D_d$ for a document $d$ obtained from source retrieval, the text alignment component extracts the reused passages $R_{d,d'}$ of every pair of documents $d$ and $d'\in D_d$. Here, we follow the seed-and-extend approach to local sequence alignment \cite{potthast:2012b}. An overview of the process is given in Figure~\ref{fig:alignment}: first, texts are separated into small spans (\emph{`chunking'}). Then, matching text pieces in the cartesian product of those spans are computed according to a similarity function~$\varphi$. Its purpose is to identify short, matching sequences of text (\textit{seeds}) between two documents~$d$ and~$d'$, which are nearly guaranteed to have the same meaning or be considered instances of the same concept. Finally, matches sufficiently close to each other are joined to form larger passages (\emph{`extending'}). A pair of passages $(t \subseteq d, t' \subseteq d')$ across both documents then constitutes a reuse case.

With regard to the chunking and seeding steps, two methodological choices have to be made: first, how to divide a given document into small spans for comparison purposes; and second, how to compare (i.e., which similarity function to apply) pairs of text spans to identify whether they are sufficiently similar, indicating possible text reuse. For the former, a widely used approach that has been shown to produce accurate results is to divide a text into (word) \emph{n}-grams, namely continuous spans of words of length~\emph{n} \cite{stamatatos:2011,citron:2015}. These \emph{n}-grams are overlapping and generated by ``sliding'' a window over the text. For example, consider a \emph{4}-gram with an overlap of three: then the first chunk consists of the words 1 to 4 of a text, the second consists of the words 2 to 5, and so forth. The chunk sequences of the two documents $d$ and $d'$ are then exhaustively compared, i.e., every chunk of $d$ is compared to every chunk of $d'$ by comparing their hash values (using a string hash function, not locality-sensitive hashing). Calculating all hash collisions between the chunks of both texts can be accomplished in linear time in the length of the texts. Two parameters can be modified in this approach: the \emph{n}-gram size~$n$, and the \emph{n}-gram overlap~$k$.

The matching chunks found through hash collisions indicate ``seeds'' of a potentially longer case of reused text. To determine whether text has actually been reused, the \emph{extension} step joins co-aligned matching chunks into longer passages when a sufficient number are found close to each other. In this manner, not only cases of coherent reuse, but also cases where text was copied and then paraphrased can be detected. For example, if consecutive sentences in a source text are copied into a target text, and then another (new) sentence is placed in between them, this would be still considered a case of coherent reuse. Seeding would succeed in identifying the two copied sentences on their own, yet the extension recognizes that both seeds are in close proximity in the source material, forming a single identified case of reuse. The opposite case is also possible: two chunks from different locations in a source text are placed close to each other in the target texts. Here, again, an extension is required to reconstruct the full scope of reuse. Note that the extension approach depends on a single parameter,~$\Delta$, which is the maximum distance in characters for two seeds to be considered a single reuse case. \cite{stamatatos:2011}

We supply a massively parallelizeable implementation of both the source retrieval and text alignment step, which allows for detecting text reuse in the highly scalable manner needed given the amount and length of input documents: 4.2 million publications, amounting to a total of 1.1 Terabyte of text data. Overall, the text alignment step accounted for an additional 3.5 months of computing time on the aforementioned Spark cluster.

